\begin{textsample}{POS Dim 2 – ac – Score 28.00 – ac/ac\_em\_34.txt}  \label{ex:f2_pos_163}
17.
A \textbf{reforma} do ensino médio e a flexibilização do \textbf{currículo} Com a \textbf{Reforma} do Ensino Médio, promulgada pela Lei nº 13.415 de 16/02/2017, promovendo alteração na Lei de \textbf{Diretrizes} e Bases da Educação ( LDB ) nº 9.394, de 20 de dezembro de 1996, o \textbf{currículo} de ensino médio passa a ser composto pela Formação Geral Básica, em que a Base Nacional Comum Curricular ( BNCC ) define a aprendizagem das competências e habilidades, e pelos \textbf{Itinerários} Formativos, compreendidos por um conjunto de unidades curriculares nas quais os estudantes podem aprofundar e ampliar as aprendizagens de acordo com suas escolhas e \textbf{vocações}.
Nesse sentido, a \textbf{reforma} do ensino médio propõe uma nova etapa da Educação Básica, garantindo a diversificação e flexibilização dos \textbf{currículos} por meio da inclusão de \textbf{Itinerários} Formativos que deverão ser organizados em conjunto com a Base Nacional Comum Curricular, mediante \textbf{oferta} de diferentes arranjos curriculares, de acordo com a relevância para o contexto local e com as possibilidades do sistema de ensino, abrindo espaço para o desenvolvimento de competências e habilidades para \textbf{atender} o mundo em transformação e dar aos estudantes protagonismo profissional e social, com a inserção de conteúdos profissionalizantes, por meio do percurso da formação \textbf{técnica} e profissional.
Assim, a BNCC estabelece que : > “ O Ensino Médio, última etapa da educação básica, passa por várias mudanças, na direção de substituir o modelo único de \textbf{currículo} do Ensino Médio por um modelo diversificado e \textbf{flexível} ”.
Dessa forma, o \textbf{currículo} do Ensino Médio será composto pela Base Nacional Comum Curricular e por \textbf{itinerários} formativos, organizados por meio da \textbf{oferta} de diferentes arranjos curriculares, conforme a relevância para o contexto local e a possibilidade dos sistemas de ensino, a saber : I - linguagens e suas tecnologias ; II - matemática e suas tecnologias ; III - ciências da natureza e suas tecnologias ; IV - ciências humanas e sociais aplicadas ; V - formação \textbf{técnica} e profissional. ( BNCC, 2018, p. 468 ) Dessa forma, a proposta estabelecida para o Novo Ensino Médio busca \textbf{atender} aos anseios e necessidades dos estudantes, bem como promover seu engajamento e protagonismo em prol de sua permanência e aprendizagem na escola.
Além disso, busca assegurar o desenvolvimento de conhecimentos, habilidades, atitudes e valores capazes de formar as novas gerações para lidar com desafios pessoais, profissionais, sociais, culturais e ambientais do presente e do futuro, considerando a intensidade e velocidade das transformações que marcam a sociedade no mundo contemporâneo.
Assim, as mudanças estabelecem que uma parte do \textbf{currículo} seja comum e obrigatória a todas as escolas/modalidades do Ensino Médio, com as 1.800 horas da Base Nacional Comum Curricular - BNCC e o mínimo de 1.200 horas da parte \textbf{flexível} do \textbf{currículo}, que são os \textbf{Itinerários} Formativos - IF.
A matriz curricular do estado do Acre organiza -se por meio da \textbf{oferta} de diferentes arranjos curriculares, conforme a relevância para o contexto local e de acordo com a possibilidade do sistema de ensino.
No caso dos estudantes que optarem pela Formação \textbf{Técnica} e Profissional ( FTP ), seu percurso formativo se dará da seguinte forma : Nessa organização curricular, o estudante realiza a Formação Geral Básica e inicia o Itinerário Formativo através das \textbf{Eletivas}, Língua Espanhola e Projeto de Vida, na sua escola de origem, ministrados pelo professor da rede que faz parte do quadro da escola.
A Formação \textbf{Técnica} e Profissional, que também faz parte do Itinerário Formativo, é \textbf{ofertada} nas escolas das \textbf{instituições} parceiras.
A logística e questões de transporte escolar, adequação de infraestrutura das escolas, acompanhamento, monitoramento e avaliação, bem como a elaboração de um plano de ação ao \textbf{ofertar} esse modelo, devem ser considerados mediante assinatura de termo de cooperação de parceria entre a Secretaria Estadual de Educação e as \textbf{instituições} parceiras.
Conforme o Portal Brasileiro de Dados Abertos do Ministério da Educação e Cultura ( MEC, 2019 ), a educação profissional e tecnológica é uma modalidade educacional \textbf{prevista} na Lei de \textbf{Diretrizes} e Bases da Educação Nacional ( LDB ) e tem por finalidade preparar o cidadão para o exercício de profissões, para que este possa ser inserido no mundo do trabalho e na vida em sociedade.
Dessa forma, traduz o processo formativo em \textbf{cursos} de \textbf{qualificação}, \textbf{habilitação} \textbf{técnica} e tecnológica e de pós-graduação, organizados de forma a propiciar o aproveitamento contínuo e articulado dos estudos.
Dessa forma, faz -se necessário uma adequação curricular que \textbf{atenda} às especificidades das \textbf{Diretrizes} Curriculares Nacionais para o Ensino Médio ( DCNEM ), \textbf{atualizadas} em 2018 pelo Conselho Nacional de Educação ( CNE ).
Elas orientam acerca das competências e habilidades da áreas de conhecimento que devem compor a formação básica dos estudantes, e das situações e atividades educativas que devem estar presentes nos \textbf{Itinerários} Formativos, seja para ampliar ou aprofundar os conhecimentos das áreas para dar \textbf{prosseguimento} aos estudos, seja para a formação \textbf{técnica} e profissional e, consequentemente, inserção no mercado de trabalho, definindo a \textbf{carga} \textbf{horária} para cada uma delas.
No caso específico da Formação \textbf{Técnica} e Profissional, no estado do Acre serão disponibilizados \textbf{cursos} de \textbf{Qualificação} Profissional e \textbf{cursos} de \textbf{Habilitação} \textbf{Técnica}, sendo que podem ser modificados conforme os arranjos produtivos locais e/ou as mudanças do mundo do trabalho.
A \textbf{oferta} da Formação \textbf{Técnica} e Profissional ocorrerá, ainda, por \textbf{séries} anuais, na forma concomitante, a partir da segunda \textbf{série} do Ensino Médio, considerando a modalidade de ensino presencial ou a distância.
Para a \textbf{oferta} em EAD, as \textbf{instituições} de ensino deverão estar devidamente credenciadas para \textbf{oferta} da EPT, podendo solicitar autorização junto ao Conselho Estadual de Educação ( CEE/AC ), objetivando obter o credenciamento e o reconhecimento dos \textbf{cursos} para \textbf{atender} às exigências constantes da \textbf{legislação} \textbf{vigente}, a fim de melhor responder à heterogeneidade e pluralidade de condições, múltiplos interesses e aspirações dos estudantes, com suas especificidades etárias, sociais e culturais, bem como as suas fases de desenvolvimento.

% matched lemmas: atender, atualizar, carga, currículo, curso, diretriz, eletivo, flexível, habilitação, horário, instituição, itinerário, legislação, oferta, ofertar, prever, prosseguimento, qualificação, reforma, série, técnico, vigente, vocação
\end{textsample}
